\section{Valuation and backtesting engine}\label{sec:engine}

The engine values a portfolio by propagating the dollar amounts held in each constituent, following the general method of \citet{milhau2026}: amounts are updated with equation~\eqref{eq:amount-update} between rebalancing dates and reset by the allocation rule on rebalancing dates.
Working with amounts rather than weights has two advantages.
First, it handles any rebalancing frequency with the same code (daily rebalancing and buy-and-hold are the two limiting cases).
Second, it gives the amount traded at each date directly, from which transaction costs are computed (Decision~D4).

\paragraph{Verification.}
The engine reproduces the worked examples of the course: the valuation of an equally weighted portfolio of five US stocks between 2 and 5 January 2015 (final value \$978.21, \$6.37 traded to return to equal weights), the drift of weights in a buy-and-hold portfolio, and the compounding of returns.
These examples are encoded as automated tests (Appendix~\ref{app:software}).

\todo{Phase 1: forward positions, transaction costs, performance and risk measures.}
